\chapter{Objetivos}

El objetivo principal de este Trabajo de Fin de Grado es lo siguiente: desarrollar un videojuego narrativo en 2D con un sistema de mecánicas adaptativas que explore una crisis de identidad desde la perspectiva de género. Surge de la intersección entre la ingeniería de software aplicada al entretenimiento interactivo y la capacidad expresiva del medio del videojuego para generar emptaía frente a realidades sociales complejas \parencite{bogost2011,navarro2023}. En consecuencia, el objetivo general del proyecto consiste en conceptualizar, diseñar y desarrollar un prototipo funcional de videojuego narrativo en \ac{2D} con perspectiva cenital oblicua (\textit{top-down} 3/4) en el motor Unity que, mediante la implementación de un sistema de mecánicas adaptativas en segundo plano y una mecánica de consentimiento explícito, permita al jugador experimentar la crisis de identidad de género del protagonista y sus consecuencias psicológicas derivadas del rechazo social e internalizado. A través de esta propuesta se busca alcanzar una armonía ludonarrativa profunda \parencite{hernandez2024}, logrando que las reglas del sistema, las habilidades del jugador y la interfaz gráfica de usuario (\ac{UI}) evolucionen de manera sincroniada con el estado anímico del protagonista sin interrumpir la jugabilidad \parencite{gomez2026}.

Para alcanzar de la mejor manera posible al objetivo final del tema planteado, se plantea los siguientes objetivos específicos:
\begin{itemize}
  \item \textbf{Realizar un estudio teórico previo y fundamentar la perspectiva psicológica del proyecto:} analizar la literatura académica sobre narrativa digital, narrativa ambiental, disonancia y consonancia ludonarrativa e identidad proyectiva \parencite{finestrat2018, hernandez2024, gomez2026}, complementando esta base documental con el asesoramiento de una especialista en psicología para garantizar un tratamiento riguroso, empático y libre de estereotipos sobre la disforia y el estrés de minorías.
  \item \textbf{Elaborar el Documento de Diseño de Juego (\ac{GDD}) completo:} definir formalmente la sinopsis argumental, la estructura de los tres mundos abstractos (\textit{Vecindad Urbana}, \textit{Laberinto Fracturado} y \textit{Azotea Libre}), las manifestaciones simbólicas del escrutinio social y el flujo de mecánicas adaptativas, estableciendo las reglas exactas que modulan la transición hacia el desenlace abierto de la obra.
  \item \textbf{Establecer la dirección artística y el diseño visual en \ac{2D} con sensación de profundidad (2.5D):} diseñar la identidad gráfica del prototipo inspirándose en referentes contemporáneos como \textit{Oxenfree} \parencite{nightschool2016}, \textit{Night in the Woods} \parencite{infinitefall2017} y \textit{Undertale} \parencite{fox2015}, especificando las paletas cromáticas evolutivas, el diseño minimalista del avatar y de los entornos, la ordenación espacial en el eje vertical (\textit{Y-Sorting}), la iluminación bidimensional (\ac{URP} 2D) y las variaciones visuales de la interfaz.
  \item \textbf{Programar un módulo evaluador del comportamiento en segundo plano:} implementar en lenguaje C\# un sistema observador desacoplado que recolecte métricas de juego de forma transparente durante la partida (tales como acciones de confrontación directa, tiempos de inmovilidad o maniobras de esquiva y rodeo), evaluando dinámicamente la tendencia conductual predominante del usuario (agresiva, pasiva o evasiva).
  \item \textbf{Desarrollar una arquitectura de software flexible y el mecanismo de Consentimiento Explícito:} construir una base de programación modular apoyada en objetos de datos (\textit{ScriptableObjects}) de Unity que permita inyectar y modificar en tiempo real tanto las habilidades del personaje como la disposición y opacidad de la \ac{UI} sin corromper el estado de la partida, integrando un panel de consulta que solicite el permiso del jugador antes de aplicar cada mutación mecánica.
  \item \textbf{Construir y validar un prototipo funcional y jugable (demo):} integrar los sistemas programados, los recursos artísticos y el diseño de niveles en un corte vertical (\textit{vertical slice}) ejecutable en PC, llevando a cabo pruebas de funcionamiento e iteración para verificar la estabilidad técnica de las transiciones adaptativas y sentar las bases para la futura expansión del videojuego tras el \ac{TFG}.
\end{itemize}
